\chapter{Produto educacional}
\label{cap:produto}

% =====================================================================
% PAPEL DESTE CAPÍTULO (exigência do PROFMAT):
%   Descreve o produto educacional como objeto apropriável por
%   professores em exercício. NÃO re-explica a arquitetura (Cap. 4)
%   nem a implementação (Cap. 5) — descreve o que o professor RECEBE.
%
%   PENDENTE: a composição final do banco curado (Seção 7.3) e o manual
%   em PDF (Seção 7.4) dependem do lote de execuções reais; as marcas
%   estão no texto.
% =====================================================================

O mestrado profissional exige que a pesquisa resulte em um produto educacional diretamente
apropriável pela prática docente. Este capítulo descreve esse produto do ponto de vista de quem
o recebe: o professor de Matemática do Ensino Médio. Os fundamentos que o justificam foram
estabelecidos nos Capítulos~\ref{cap:contexto} a~\ref{cap:didatica}; seu projeto e construção,
nos Capítulos~\ref{cap:arquitetura} e~\ref{cap:implementacao}. Aqui, a pergunta é outra: o que,
concretamente, um professor em exercício pode obter, instalar e usar --- e em que condições.

\section{Visão geral: três entregáveis articulados}
\label{sec:produto-visao}

O produto educacional é composto por três entregáveis que se complementam, todos versionados em
repositório público:

\begin{enumerate}
\item O \textbf{sistema multiagente} --- o \emph{software} que gera, verifica e avalia questões
  a partir dos parâmetros do professor, com interface web local (Seção~\ref{sec:produto-sistema}).
\item O \textbf{banco curado de questões} --- um conjunto de questões aprovadas pelo ciclo
  completo do sistema, categorizadas por tema, habilidade BNCC, nível cognitivo e dificuldade,
  utilizável mesmo por quem não instalar o sistema (Seção~\ref{sec:produto-banco}).
\item O \textbf{manual do professor} --- o documento que ensina a instalar, usar e, sobretudo,
  a compreender os limites da ferramenta (Seção~\ref{sec:produto-manual}).
\end{enumerate}

A articulação é deliberada: o banco atende ao professor que quer \emph{material pronto}; o
sistema atende ao que quer \emph{produzir o próprio material}; o manual conecta os dois e
sustenta o uso crítico que o Capítulo~\ref{cap:contexto} argumentou ser indelegável.

\section{O sistema multiagente}
\label{sec:produto-sistema}

O professor recebe um pacote Python instalável com uma interface web que roda localmente no
próprio computador --- nenhum dado da sua prática é enviado a servidores do projeto. A interface
organiza o trabalho em duas áreas. Em \emph{gerar}, o professor preenche o formulário de
especificação, indica quantas questões deseja e acompanha cada uma chegando com a
\textbf{trilha do ciclo}: o que o Gerador redigiu, o que o Verificador Simbólico recalculou de
forma independente e como o Crítico Didático pontuou. Em \emph{banco curado}, consulta o acervo
com filtros, registra sua própria avaliação de cada questão e seleciona itens para exportar a
lista em Markdown, \LaTeX{} ou Word, nas versões do aluno e do professor.

A decisão de exibir o resultado recalculado pelo verificador --- e não apenas um selo de
``aprovado'' --- é deliberada. O erro silencioso descrito na Seção~\ref{sec:erro-silencioso} é
justamente aquele que se apresenta com aparência de correção; mostrar ao professor \emph{o que}
a máquina calculou por conta própria devolve a ele a possibilidade de discordar, em vez de
substituir a confiança cega no LLM por confiança cega no sistema.

\textbf{Registro da avaliação do professor.} Cada questão do banco admite um julgamento humano
--- aceita, aceita com ajuste ou recusada --- acompanhado de um comentário livre. O registro
serve a dois propósitos: organiza a curadoria do próprio professor e constitui a matéria-prima
da avaliação empírica relatada no Capítulo~\ref{cap:avaliacao}, ao permitir contrastar a
aprovação pelos agentes com o juízo de quem conhece a turma.

\textbf{Depósito em pasta sincronizada.} O professor pode indicar uma pasta que um serviço de
nuvem (Google Drive, OneDrive) já sincronize. A cada questão salva, o sistema grava ali um
arquivo legível com a questão e sua verificação, um arquivo com o ciclo completo e uma planilha
com uma linha por questão --- veredicto do verificador, nota mínima do crítico, número de
iterações e avaliação do professor. Não há autenticação nem chamada de rede: quem sincroniza é
o cliente da nuvem já instalado, escolha que evita depender de credenciais que expiram e mantém
o princípio de que nenhum dado transita por servidores do projeto.

\textbf{Requisitos mínimos.} Um computador com Python 3.11 ou superior e acesso a um provedor
de modelo de linguagem. Este último admite dois caminhos, e a diferença importa para a
realidade da escola pública:

\begin{itemize}
\item \textbf{Com API comercial} (Anthropic ou OpenAI): exige uma chave de API paga pelo uso;
  o custo por questão é pequeno, mas existe, e a chave é informada diretamente no painel de
  configuração da interface --- não fica gravada em arquivo algum.
\item \textbf{Sem API paga} (Ollama): o professor instala um modelo aberto que roda no próprio
  computador, sem chave e sem custo por uso. O requisito desloca-se do financeiro para o de
  \emph{hardware} (memória suficiente para o modelo escolhido), e a qualidade da geração é
  inferior à dos modelos comerciais --- limitação documentada no manual.
\end{itemize}

\textbf{Caminhos de instalação.} Para o professor com alguma familiaridade com linha de
comando, a instalação são três comandos, reproduzidos no manual e no repositório. Para o
professor sem essa familiaridade, o manual oferece um roteiro passo a passo com capturas de
tela --- e, de forma realista, reconhece que o caminho de menor atrito para esse perfil é
começar pelo banco curado, que não exige instalação alguma.

\section{O banco curado de questões}
\label{sec:produto-banco}

O banco reúne as questões que atravessaram o ciclo completo --- verificação simbólica e
avaliação didática --- e foram aprovadas. Cada questão carrega seus metadados (tema, habilidade
BNCC, nível cognitivo, dificuldade, natureza e formato) e o histórico do ciclo que a produziu,
incluindo o veredicto do verificador. É essa categorização que permite ao professor localizar
rapidamente, por exemplo, questões \emph{aplicadas} de função exponencial no nível
\emph{analisar} associadas à habilidade \texttt{EM13MAT304}.

O banco é distribuído em dois formatos: \textbf{JSON} (o formato nativo do sistema, que
preserva todos os metadados e permite reimportação e filtragem programática) e \textbf{PDF}
(listas organizadas por tema, com gabarito e resolução comentada, para uso direto sem qualquer
ferramenta).

% PENDENTE: composição final do banco (meta do projeto: 60-100 questões
% aprovadas, equilibradas por tema e nível) — depende do lote de
% execuções reais com API. Atualizar os números abaixo após a curadoria.
\emph{[A composição quantitativa final do banco --- meta de 60 a 100 questões aprovadas,
equilibradas por tema e nível cognitivo --- será consolidada após o lote de geração com
provedor real e a curadoria final, e reportada nesta seção.]}

\textbf{Licenciamento.} O banco é distribuído sob licença Creative Commons
Atribuição-CompartilhaIgual (CC BY-SA), que autoriza o professor a usar, adaptar e redistribuir
as questões --- inclusive modificando enunciados para sua turma --- desde que mantenha a
atribuição e a mesma licença. O código do sistema mantém licença de \emph{software} livre
própria, indicada no repositório.

\section{O manual do professor}
\label{sec:produto-manual}

O manual é o entregável que transforma o \emph{software} em produto educacional: sem ele, o
sistema é acessível apenas a quem já domina o vocabulário técnico. Sua estrutura tem quatro
partes:

\begin{enumerate}
\item \textbf{Instalação} --- os dois caminhos (com e sem API paga), em roteiro passo a passo.
\item \textbf{Uso típico} --- como preencher cada parâmetro da especificação, com exemplos
  comentados de boas especificações e do efeito de cada parâmetro sobre a questão gerada.
\item \textbf{Fluxos de trabalho recomendados} --- três roteiros práticos: gerar questões
  avulsas para uma avaliação; construir um acervo pessoal ao longo do bimestre; e montar listas
  diferenciadas por nível a partir do banco.
\item \textbf{Limitações conhecidas e como contorná-las} --- a parte mais importante do
  documento. Explica, em linguagem não técnica, o que significa o veredicto \emph{não-verificável}
  (e por que essas questões pedem conferência manual prioritária), por que a qualidade varia
  com o provedor escolhido, e por que a validação final de \emph{toda} questão permanece sendo
  atribuição do professor --- retomando, agora em registro prático, o argumento da
  Seção~\ref{sec:erro-silencioso}.
\end{enumerate}

% PENDENTE: o manual em PDF será derivado da documentação do repositório
% (README do sistema) após a estabilização da interface e do banco.
\emph{[O manual em PDF encontra-se em elaboração a partir da documentação do repositório; a
versão final acompanha o produto no repositório público e no Apêndice desta dissertação.]}

\section{Estratégia de disseminação}
\label{sec:produto-disseminacao}

A disseminação apoia-se em três frentes. A primeira é o \textbf{repositório público}: código,
banco e manual versionados juntos, com histórico aberto --- qualquer professor, ou pesquisador,
pode obter a versão mais recente, relatar problemas e propor melhorias. A segunda é a
\textbf{rede do PROFMAT}: apresentação do produto em encontros do programa e disponibilização
às turmas de egressos, cujo perfil --- professores em exercício com formação continuada em
Matemática --- coincide exatamente com o usuário-alvo. A terceira são as \textbf{parcerias com
escolas}: o desenho da avaliação por painel docente (Capítulo~\ref{cap:avaliacao}) já estabelece
contato com professores da rede, e esses mesmos professores são os multiplicadores naturais do
produto em suas escolas.

Cabe, por fim, registrar o que o produto \emph{não} pretende ser. Ele não substitui o livro
didático, não substitui bancos institucionais consolidados como os do ENEM e da OBMEP e ---
sobretudo --- não substitui o julgamento profissional docente. Ele ataca um ponto específico e
bem delimitado da prática: o custo, em tempo e em risco de erro, de produzir questões novas,
alinhadas à BNCC e calibradas por nível cognitivo. O que o sistema automatiza é a produção e a
primeira camada de verificação; a autoria pedagógica --- decidir o que avaliar, quando e para
quem --- permanece onde sempre esteve.
